\section{Pr\'actica 3: LEDs en cascada}

\subsection{Introducción}
\noindent
En esta práctica, se colocaron 24 leds en 8 filas de 3 leds cada una en una protoboard; 
el objetivo es que los 24 leds formarán 3 secuencias distintas; la primera era que cada 
fila se apagará y prendiera empezando por las filas de los extremos (0 y 7), terminando 
en las filas del centro (3 y 4) y luego de regreso, la segunda secuencia era que empezaran 
a prenderse y apagarse una fila por una, empezando en la fila 0 y terminando en la fila 
7, por último, la tercera secuencia inicia con la fila 0 encendiéndose, luego la fila 1 
pero sin apagarse la fila 0 y así sucesivamente hasta que todas las filas estén encendidas.

\subsection{Materiales}
\begin{itemize}
	\item PIC16F628A
	\item 8 resistencias de 220$\Omega$
	\item Alambre para Protoboard
	\item 1 Cristal de cuarzo de 4MHz
	\item 2 Capacitores de 22pf
	\item Fuente de poder
\end{itemize}

\subsection{Software}
\begin{itemize}
	\item Proteus 8.0
	\item MPLAB X IDE
	\item PICkit 3 Programmer App and Scripting Tool v3.10
\end{itemize}

\subsection{Circuito}
\begin{figure}[H]
	\centering
	\includegraphics[width = 0.8 \textwidth]{PROY_03.png}
	\caption{Circuito esportado de Proteus}
	\label{fig:proy03}
\end{figure}

\subsection{Desarrollo}
\noindent
Comenzamos grabando el pic con el archivo .hex (previamente probado en la simulación de 
proteus) haciendo uso de nuestro aparato Pickit3, entonces basándonos en el circuito 
hecho en Proteus procedimos a montarlo en la protoboard; sin embargo, nos percatamos que 
este mismo no prendio a la primera por lo que fue necesario checar las conexiones 
(Figura \ref{fig:conex03}) y procedimos a volver a acomodar los leds y los alambres de 
la protoboard para que el circuito se viera más ordenado y asegurarnos de que todo 
estuviera en orden (Figura \ref{fig:leds03} y Figura \ref{fig:circuitterminado03}).

\begin{figure}[ht]
	\centering
	\includegraphics[width = 0.5 \textwidth]{conex_03.png}
	\caption{Revisando las conexiones}
	\label{fig:conex03}
\end{figure}

\begin{figure}[ht]
	\centering
	\includegraphics[width = 0.5 \textwidth]{leds_03}
	\caption{Acomodando los leds en la protoboard}
	\label{fig:leds03}
\end{figure}

\begin{figure}[ht]
	\centering
	\includegraphics[width = 0.5 \textwidth]{reacomodo_03.png}
	\caption{Re-acomodando los leds en la protoboard}
\end{figure}

\noindent
Después procedimos a terminar el armado y verificar nuevamente que todo estuviera en su 
lugar, descubriendo que el error cometido anteriormente fue que habían alambres mal 
conectados, mismo error que fue corregido más adelante.

\begin{figure}[ht]
	\centering
	\includegraphics[width = 0.5 \textwidth]{circuito_terminado_03.png}
	\caption{Ciruito terminado}
	\label{fig:circuitterminado103}
\end{figure}

\noindent
Una vez corregidos los errores procedimos a probar el circuito haciendo uso de nuestra 
fuente de poder, el resultado fue satisfactorio pues las secuencias programadas se 
concretaron como estaban programadas.

\begin{figure}[ht]
	\centering
	\includegraphics[width = 0.5 \textwidth]{circuito_funcionando1_03.png}
	\caption{Ciruito en funcionamiento}
	\label{fig:circuitfun03}
\end{figure}

\begin{figure}[ht]
	\centering
	\includegraphics[width = 0.5 \textwidth]{circuito_funcionando2_03.png}
	\caption{Ciruito terminado}
	\label{fig:circuitterminado203}
\end{figure}
\clearpage

\subsection{C\'odigo en C}
\begin{verbatim}
// PIC16F628A Configuration Bit Settings

// 'C' source line config statements

// CONFIG
#pragma config FOSC = XT        // Oscillator Selection bits 
//(XT oscillator: Crystal/resonator on RA6/OSC2/CLKOUT and RA7/OSC1/CLKIN)
#pragma config WDTE = OFF       // Watchdog Timer Enable bit (WDT disabled)
#pragma config PWRTE = ON       // Power-up Timer Enable bit (PWRT enabled)
#pragma config MCLRE = ON       // RA5/MCLR/VPP Pin Function Select bit 
//(RA5/MCLR/VPP pin function is MCLR)
#pragma config BOREN = ON       // Brown-out Detect Enable bit (BOD enabled)
#pragma config LVP = OFF        // Low-Voltage Programming Enable bit 
//(RB4/PGM pin has digital I/O function, HV on MCLR must be used for programming)
#pragma config CPD = OFF        // Data EE Memory Code Protection bit 
//(Data memory code protection off)
#pragma config CP = OFF         // Flash Program Memory Code Protection bit 
//(Code protection off)

// Use project enums instead of #define for ON and OFF.

#include <xc.h> // Librería principal del microcontrolador
#define _XTAL_FREQ 4000000 // Frecuencia del cristal = 4 MHz 

// Arreglos para cada secuencia

unsigned char secuencia12[8] = {0x01,0x02,0x04,0x08,
				0x02,0x08,0x20,0x80}; // Secuencias 1 y 2 
unsigned char secuencia3[8] = {0x01,0x03,0x07,0x0F,
				0x02,0x0A,0x2A,0xAA}; // Secuencia 3

void main (void){

TRISA = 0b11100000; // RA0–RA4 como salidas, RA5–RA7 como entradas
TRISB = 0b00000000; // Todo PORTB como salida
 
// Limpieza de los puertos
PORTA = 0x00; 
PORTB = 0x00; 

while(1){ 

//Secuencia 1

// Empieza de afuera hacia adentro
for(int i = 0; i<4; i++){ 
PORTA = secuencia12[i]; 
PORTB = secuencia12[7-i]; 
__delay_ms(400); 

}

// De adentro hacia afuera
for(int i = 2; i>=0; i--){ 
PORTA = secuencia12[i]; 
PORTB = secuencia12[7-i]; 
__delay_ms(400); 

}

// Limpieza de los puertos 
PORTA = 0x00; 
PORTB = 0x00; 
__delay_ms(400); 


// Secuencia 2

// Los leds prenden en orden
for(int i = 0; i<4; i++){ 
PORTA = secuencia12[i]; 
PORTB = 0x00; 
__delay_ms(400);
}

for(int i = 4; i<8; i++){ 
PORTA = 0x00; 
PORTB = secuencia12[i]; 
__delay_ms(400);
}

// Limpieza de los puertos 
PORTA = 0x00; 
PORTB = 0x00;
__delay_ms(400);


// Secuencia 3

// Se prenden los leds del puerto A
for(int i = 0; i<4; i++){ 
PORTA = secuencia3[i]; 
PORTB = 0x00; 
__delay_ms(400);
}

// Se prenden los leds del puerto B sin apagar los del puerto A
for(int i = 4; i<8; i++){ 
PORTA = 0x0F; 
PORTB = secuencia3[i]; 
__delay_ms(400);
}

\end{verbatim}

\subsection{Circuito cuando se ejecuta la simulación}
\begin{figure}[H]
	\centering
	\includegraphics[width = 0.5 \textwidth]{PROY_03_sim.png}
	\caption{Ciruito exportado de Proteus}
	\label{fig:proy03sim}
\end{figure}

\subsection{Conclusi\'on}
\noindent
Logramos desarrollar y programar tres secuencias de LEDs usando el PIC16F628A, 2 
arreglos (uno para las secuencias 1 y 2, el otro para la secuencia 3) y ciclos for, 
lo cual ayudó a organizar mejor el código. También se reforzó el uso del simulador para 
verificar el funcionamiento del programa antes de implementarlo físicamente, pues en la 
simulación en proteus el programa funcionó bien y las secuencias de LEDs se ejecutaron 
sin errores, pero en el circuito no se logró el funcionamiento correcto al primer 
intento, por lo que tuvimos revisar cuidadosamente las conexiones y cambiar algunas 
conexiones. 
